\section{Classification and interior regularity}
\label{sec:assembly}

We now pass from the exclusion of compact pieces to the local geometry of the
original singular set.  The only limiting sets still available will be the
three models in the epsilon-regularity theorem.

\begin{corollary}[Classification of the generalized limits]
\label{cor:classification}
The closed set of every global generalized minimizer arising from absolute
minimizers is empty, a line, a $120$-degree triod of half-lines, or a half-line.
\end{corollary}

\begin{proof}
Let $(v,H,\{p_{kl}\})$ be such a minimizer.  If $H$ disconnects the plane,
Proposition~\ref{prop:planar-facts}(\ref{item:classification}) gives the
conclusion.  Otherwise Proposition~\ref{prop:planar-facts}(\ref{item:connected-absolute})
gives absolute minimality.  If $H$ had a bounded component,
Lemma~\ref{lem:compact-isolation} would enclose it in a separated compact
piece $A$ of positive length.  Its length measure is concentrated on points
with arc neighborhoods, by
Proposition~\ref{prop:planar-facts}(\ref{item:global-estimates}), so one such
arc lies in $A$.  Proposition~\ref{prop:compact-piece} rules out this piece.
Thus $H$ has no bounded component, while
Proposition~\ref{prop:planar-facts}(\ref{item:classification}) allows at most
one unbounded component.  Hence $H$ is empty or connected.  In either case all but
at most one component lie in a compact set, so the remaining classification
statement in that proposition applies.
\end{proof}

\begin{proof}[Completion of the proof of Theorem~\ref{thm:main}]
Proposition~\ref{prop:initial-reduction} allows us to apply
Proposition~\ref{prop:planar-facts}.  Fix $x\in K$ and take a blow-up with
limiting set $H$.  Since $0\in H$, Corollary~\ref{cor:classification} leaves
three possibilities: a line, a half-line, or a triod of half-lines.
Choose $\rho>0$ so that in $B_{3\rho}(0)$ the limiting set is a model centered
at zero.  If the vertex of a half-line or triod is different from zero, it
suffices to take $3\rho$ less than its distance from zero; the model is then a
line.

Local bilateral Hausdorff convergence gives
\[
 \operatorname{dist}_{\mathrm H}
 \bigl(K_j\cap B_{2\rho}(0),H\cap B_{2\rho}(0)\bigr)=o(\rho).
\]
To justify truncation at the sphere, move each model point near
$\partial B_{2\rho}$ slightly inward along its ray, apply convergence in
$B_{3\rho}$, and then let the inward displacement tend to zero.  Approximating
points are now inside $B_{2\rho}$; the reverse distance follows by the same
radial projection onto the truncated model.

Put $s_j=r_j\rho$ and return to the original coordinates.  For large $j$, the
closed set in $B_{2s_j}(x)$ has model distance $o(s_j)$, the enlarged ball is
compactly contained in the original domain, and $s_j\leq1$.  The fixed-data
epsilon-regularity statement
Proposition~\ref{prop:planar-facts}(\ref{item:epsilon}) now applies to the
original minimizer.  It follows that the entire set $K\cap B_{s_j}(x)$ is
$C^{1,\alpha}$-diffeomorphic to the corresponding model.

The endpoint or triple junction in this chart may have shifted from its
center.  This causes no difficulty: $x$ is either an interior point of one of
the arcs or the actual endpoint or junction.  Restricting the chart around
$x$ gives its corresponding local model.  The equal-angle conclusion for a
triple junction is part of the epsilon-regularity theorem.  We have therefore
proved the asserted interior regularity at every point of $K$.

It remains to prove local finiteness of the global connected components.
For each $x\in K$, retain one of the original epsilon-regularity disks
$B_{s_x}(x)\Subset\Omega$, whose entire intersection with $K$ is connected.
If $U\Subset\Omega$ is open, then
$K\cap\overline U$ is compact, so finitely many of these disks cover it.
Each disk meets only one global connected component of $K$, because its entire
intersection with $K$ is connected.  Every global component meeting $U$
therefore belongs to the finite list represented by the covering disks.
\end{proof}

\begin{remark}[Local scope]
The local models also make the endpoints and triple junctions a relatively
closed discrete subset of $K$, so only finitely many of them meet a compact
subset of $\Omega$.  Define the edges to be the connected components left
after deleting these endpoints and triple junctions from $K$.  In each model
chart, deleting those points leaves at most three connected pieces, each
contained in one global edge.  A finite model-chart cover therefore meets only
finitely many global edges.  These are interior finiteness statements, with
no assertion at $\partial\Omega$.  They concern global components rather than
components of $K\cap U$: even one straight line can have infinitely many
intersection components with an arbitrary relatively compact open set.
\end{remark}

The geometric conclusion also gives an endpoint integrability estimate for
the gradient.  Its exponent is suggested by the crack-tip profile
$r^{1/2}\sin(\theta/2)$: the gradient has size proportional to $r^{-1/2}$,
so its superlevel sets near the tip have area of order $t^{-4}$ for large
$t$, whereas its fourth power is not locally integrable
\cite[Section~1.4]{delellis_focardi2025}.  The next result gives this
weak-$L^4$ bound for every minimizer in Theorem~\ref{thm:main}, using the
De Lellis--Focardi equivalence between the absence of irregular points and
the weak-$L^4$ estimate.

\begin{corollary}[Local weak-$L^4$ gradient bound]\label{cor:weak-l4}
Let $(u,K)$ satisfy the hypotheses of Theorem~\ref{thm:main}.
Let $\nabla u$ denote the approximate gradient furnished by
Proposition~\ref{prop:initial-reduction}; it agrees almost everywhere on
$\Omega\setminus K$ with the gradient in \eqref{eq:original-energy}.
For every open $U\Subset\Omega$ there is a constant $C_U<\infty$,
depending on $U$ and the minimizer, such that
\begin{equation}\label{eq:weak-l4}
 \mathcal L^2\bigl(\{x\in U\setminus K:|\nabla u(x)|>t\}\bigr)
 \le C_Ut^{-4}\qquad\text{for every }t>0,
\end{equation}
where $\mathcal L^2$ denotes planar Lebesgue measure. In particular,
$\nabla u\in L^{4,\infty}_{\mathrm{loc}}(\Omega)$ and
$\nabla u\in L^p_{\mathrm{loc}}(\Omega)$ for every $0<p<4$.
\end{corollary}

\begin{proof}
Fix $U\Subset\Omega$ and choose a smooth open set $W$ with
$U\Subset W\Subset\Omega$. Proposition~\ref{prop:initial-reduction}
makes the restricted pair a reduced rectifiable absolute minimizer for
the energy $E_1$ of De Lellis--Focardi, with bounded fidelity datum.
The three models in Theorem~\ref{thm:main} are precisely their regular
pure-jump, loose-end, and triple-junction models: the arcs can be
straightened to their tangent rays by a $C^1$ chart tangent to the
identity at the center, and the triple angles are $120$ degrees
\cite[Definitions~1.3.2 and~1.5.1]{delellis_focardi2025}.
Thus the irregular stratum $K^{(i)}$ of the restricted pair is empty.
Their weak-$L^4$ equivalence
\cite[Theorem~1.5.4; see also Theorem~6.1.6]{delellis_focardi2025}
gives the superlevel estimate \eqref{eq:weak-l4} on $U$, after increasing
$C_U$ for $0<t<1$. Since $\mathcal H^1(K)<\infty$ implies
$\mathcal L^2(K)=0$, this is also the weak-$L^4$ bound for the approximate
gradient. For $0<p<4$, the layer-cake formula bounds
$\int_U|\nabla u|^p\,dx$ by
\[
 p\int_0^\infty t^{p-1}
   \min\{\mathcal L^2(U),C_Ut^{-4}\}\,dt<\infty.
\]
\end{proof}
